This paper asks how love fares under power. Across the \emph{Philosophy of Intimacy and the Theory of Justice} series, intimate relation has been theorised as a generative-relational process---a coupling of two subjects out of which value is generated rather than extracted. But that account was developed, in effect, in a clearing: it did not confront the plain fact that intimate relation unfolds in a field saturated with asymmetries, which can deform, capture, freeze, and on occasion shatter the very generativity the series describes. The present paper confronts power directly.

It is organised around three questions, each answered by a different body of theory and all bearing on a single object---\emph{love}, in the precise sense of the dynamics of the real between two subjects, the part of their bond that circles the unsymbolisable and cannot be counted, traded, or directly managed. \emb{How does love create a world?}---the question of the generative-relational being (GRB), whose affirmative findings the paper takes as established, adding an \emph{etiology} of how the world-creating coupling breaks down into estrangement (\zh{すれ違い}, \emph{surechigai}). \emb{How is love not destroyed by power?}---the question for which the paper reconstructs the theory of international relations, stripping away the realist assumption that the currency of an anarchic field is relative power, and recovering IR's deeper object: how opaque subjects, unable to enter one another directly, build trust and order under uncertainty and without a central authority. \emb{How does love become free within history?}---the question for a dialectical historical materialism that holds the determining force of material conditions together with the relative autonomy of the symbolic and affective life they shape.

Three questions with an ascending logic---\emb{creation, preservation, liberation}---and a single methodological commitment: that none of the three theories can touch their shared object directly. The intimate relation is a two-level system. At the level of the \emph{end} stands the dynamics of love, a process of the real that cannot be measured or exchanged; at the level of the \emph{means} stands the symbolic order of power, calculable and strategic. The means serve the end but tend intrinsically to damage it. Hence the paper's central claim: the tools of power and wealth can protect the \emph{conditions} under which love occurs, but cannot produce love itself---a \emph{political economy of conditions}, never of love. Trust is the interface that crosses both levels, and the variable that decides whether an asymmetry of power becomes generative or frozen.

Two further claims follow. Estrangement is, in the first instance, not a malfunction but a necessary moment of the bond's self-maintenance; a relation at perfect concordance would be a relation at equilibrium, which is to say dead. And power can neither eliminate itself nor be eliminated entirely; the only viable course is to generate, alongside good value, a continual counter-force that keeps power from solidifying. \emb{A generative relation does not oppose power; it opposes the freezing of power.} This, the paper argues at its summit, is not an external ethical demand but the political appearance of an ontological law generativity already obeys---the law by which desire, generating its symbols, withholds itself from total symbolisation in order to remain alive.

This is \pinkref{Paper~XVII} of the series. It extends the account of trust opened in \pinkref{Paper~XIII} \parencite{paperxiii}, the holonomic criterion of \pinkref{Paper~IX} \parencite{paperix}, and the political economy of \pinkref{Papers~XIV--XVI}, and draws throughout on the framework of the mother paper \parencite{wan2024grb}.
